\documentclass[runningheads]{llncs}

\usepackage[T1]{fontenc}
\usepackage[utf8]{inputenc}
\usepackage[spanish,es-tabla,es-noshorthands]{babel}
\usepackage{booktabs}
\usepackage{xcolor}
\usepackage{hyperref}

% llncs solo traduce "Keywords" al francés y al alemán
\renewcommand{\keywordname}{{\bfseries Palabras clave:}}

% Texto pendiente de los resultados del CESGA: buscar \pendiente antes de entregar
\newcommand{\pendiente}[1]{\textcolor{red}{[#1]}}

\begin{document}

\title{¿Humano o máquina? Detección distribuida de texto generado por IA con Spark MLlib}
\titlerunning{Detección de texto generado por IA con Spark MLlib}

\author{Alicia Losada Sánchez \and Carlos Francisco Orejuela Lluncor \and Hugo Varela Sánchez}
\authorrunning{A. Losada, C. F. Orejuela y H. Varela}
\institute{Aprendizaje Automático a Gran Escala, Universidade da Coruña}

\maketitle

\begin{abstract}
Clasificamos textos en inglés como escritos por personas o generados por modelos de lenguaje,
usando 1,39 millones de textos y Spark MLlib. Antes de entrenar eliminamos los atajos que deja
la construcción del dataset (longitud y marcas de formato propias de cada fuente). Comparamos
cuatro familias de features (léxicas, de caracteres, estilométricas y de previsibilidad según
un modelo de n-gramas), varios clasificadores y su combinación por stacking. El mejor modelo
individual, una regresión logística sobre unigramas y bigramas, alcanza una AUC-PR de 0,837 en
test frente a una línea base de 0,225. El análisis por fuente muestra que parte de ese acierto
viene del género del texto y no de la autoría: entre fuentes del mismo dominio el problema es
bastante más difícil.

\keywords{Detección de texto generado \and Spark MLlib \and Estilometría \and Stacking}
\end{abstract}

% ------------------------------------------------------------------------------------------
\section{Introducción}

Los modelos de lenguaje producen texto difícil de distinguir del humano, lo que afecta a la
evaluación académica, a la moderación de contenidos y a la fiabilidad de lo que se publica en
la web. Los detectores más precisos usan otro modelo de lenguaje grande para medir la
probabilidad del texto~\cite{gehrmann2019gltr,mitchell2023detectgpt}, lo que es caro de aplicar
a millones de documentos. Nos preguntamos hasta dónde llegan features clásicas, baratas y
calculables en paralelo con Spark~\cite{zaharia2016spark,meng2016mllib}.

Usamos \emph{AI Text Detection Pile}: 1.392.522 textos (1,9~GB en 7 ficheros
Parquet), un 73,8\,\% humanos y un 26,2\,\% generados por IA. Es lo bastante grande para que
el preprocesado, la extracción de features y la búsqueda de hiperparámetros no quepan
cómodamente en una sola máquina, y por eso la ejecución final se hizo en el CESGA.

% ------------------------------------------------------------------------------------------
\section{Datos y preprocesado}

\paragraph{Fuentes.} El dataset solo indica si cada texto es humano o IA, pero los \texttt{id}
están ordenados por origen. Analizando la longitud y las marcas de formato a lo largo del
\texttt{id} identificamos ocho bloques que coinciden en tamaño con lo que declara la ficha:
ensayos de IvyPanda, WebText, respuestas humanas de HC3 y relatos de Reddit
WritingPrompt en el lado humano; GPT-3, GPT-2,
GPT-J y ChatGPT en el de la IA. Esta columna (\texttt{subsource}) es inferida y solo se usa
para analizar resultados, nunca como feature, porque el \texttt{id} separa las clases al 100\,\%.

\paragraph{Atajos.} Cada fuente deja huellas que no tienen que ver con quién escribe. La más
fuerte es la longitud: la mediana es de 305 palabras en la IA y 190 en los humanos, porque
los relatos de Reddit vienen recortados a unos 1.000 caracteres y GPT-2 llega a su límite de
generación. También hay marcas de formato casi exclusivas de una clase: comillas de \LaTeX{}
en el 52\,\% de Reddit, puntuación separada por espacios en el 85\,\% de HC3, el título
``\ldots{} Essay'' en el 56\,\% de IvyPanda, un espacio inicial en el 98\,\% de GPT-3 o doble
espacio tras el punto en el 60\,\% de GPT-J.

\paragraph{Limpieza.} Aplicamos las mismas reglas a las dos clases, sin mirar la etiqueta y
solo con funciones nativas de Spark (sin UDF de Python): comillas normalizadas, títulos y
asteriscos eliminados, puntuación separada unida de nuevo y espacios normalizados. Además
descartamos siempre la primera frase. No hay forma fiable de saber si un trozo empieza a mitad
de frase (puede empezar con mayúscula), y así todos los textos arrancan a mitad de documento,
como los de Reddit. El precio es perder aperturas que sí son estilo, como el ``It's important
to understand\ldots'' de ChatGPT. Tras la limpieza todas las marcas medidas bajan a cerca del
0\,\%.

\paragraph{Ventana y duplicados} Cada texto se recorta a sus primeras 100
palabras y se descartan los más cortos. Es el compromiso entre conservar Reddit (sus trozos
rondan las 170 palabras) y tener texto suficiente para la estilometría. Quedan unos 1,19
millones de textos con un 22\,\% de IA, aunque las fuentes de respuestas cortas pierden mucho
(sobrevive el 14,5\,\% de GPT-3 y el 25\,\% de HC3). Después eliminamos los duplicados exactos,
conservando la copia de menor \texttt{id}, y los textos que aparecen con las dos etiquetas.

% ------------------------------------------------------------------------------------------
\section{Features}

\begin{enumerate}

    \item \textbf{Léxicas.} Unigramas y bigramas de palabras en minúscula, con y sin
\emph{stopwords}, contados con \texttt{HashingTF} ($2^{18}$ posiciones).

\item \textbf{Caracteres.} N-gramas de 3 a 5 caracteres, conservando mayúsculas, espacios y
puntuación. MLlib no tiene n-gramas de caracteres, así que los construimos con
\texttt{RegexTokenizer} (un token por carácter) seguido de \texttt{NGram}.

\item \textbf{Estilométricas.} 32 medidas por texto: riqueza
léxica (TTR, hápax), histograma y media de la longitud de palabra, número y longitud de las
frases y su variabilidad, frecuencia de ocho signos de puntuación, mayúsculas, dígitos,
contracciones, repetición de bigramas y trigramas y proporción de \emph{stopwords}.

\item \textbf{Previsibilidad.} Un modelo de trigramas con \emph{Stupid
Backoff} entrenado con los textos humanos de \texttt{lm\_fit}. Los
conteos son un \emph{word count} distribuido. La primera versión lo hacía con RDDs
(\texttt{flatMap} y \texttt{reduceByKey}); la sustituimos por \texttt{explode} y
\texttt{groupBy} sobre DataFrames, porque el RDD serializa cada n-grama entre la JVM y Python,
mientras que con DataFrames todo el conteo se queda en la JVM. Para puntuar cada palabra, las tablas de
n-gramas se cruzan con los textos mediante \emph{broadcast joins} sobre claves hash. De cada
texto salen 9 features: media, desviación y percentil 10 de la log-probabilidad, nivel del
backoff usado y, siguiendo a GLTR~\cite{gehrmann2019gltr}, la fracción de palabras que estaban
entre las 1, 10 y 100 continuaciones más probables.
\end{enumerate}

% ------------------------------------------------------------------------------------------
\section{Modelos y metodología}

Sobre las features dispersas (familias 1 y 2) entrenamos Naive Bayes multinomial con los
conteos, y regresión logística y \texttt{LinearSVC} con TF-IDF. Sobre las densas (3 y 4)
usamos regresión logística con \texttt{StandardScaler}, GBT y Random Forest.

Con respecto a la partición de los datos: 70\% destinado a entrenamiento, 15\% a validación y
el restante a test. Aunque cabe aclarar que el 30\% de los textos humanos fueron reservados para
construir el modelo de n-gramas de la familia 4 (\texttt{lm-fit}).

Los hiperparámetros se eligen con \texttt{CrossValidator} de 3 folds sobre una muestra
estratificada del 15\,\% de train, con folds también estratificados. Las rejillas son pequeñas
(de 3 a 6 combinaciones) porque cada una cuesta tres ajustes. Con la mejor combinación se
reentrena el \texttt{Pipeline} con todo train.

\paragraph{Métricas.} La métrica principal es la AUC-PR de la clase IA. Con un 22,5\,\% de
positivos la AUC-ROC es optimista, porque los muchos verdaderos negativos la inflan, mientras
que la AUC-PR se centra en la clase que nos interesa. Damos además la AUC-ROC y la tasa de
detección con un 1\,\% de falsos positivos, que es lo relevante si acusar a un humano tiene
un coste alto. Para las métricas con umbral (precisión, recall, F1, FPR) el umbral se fija en
val maximizando el F1 y se aplica sin tocar a test. Por último, calculamos el recall en cada
generador y el FPR en cada fuente humana.

\paragraph{Stacking.} Tomamos el mejor modelo de cada familia según su AUC-PR en
val. Un meta-modelo (escalado y regresión logística) aprende a
combinar sus scores en val, donde los modelos base no han visto los textos; en train sus
scores serían optimistas. Medimos lo que aporta cada familia reentrenando el meta-modelo sin
ella (ablación).

% ------------------------------------------------------------------------------------------
\section{Organización del código}

El código está en \texttt{src/} y \texttt{main.py} lo ejecuta de principio a fin:
\texttt{clean\_data.py} (carga, limpieza, ventana, duplicados y particiones),
\texttt{features\_families/} (un módulo por familia, cada uno escribe su Parquet),
\texttt{models.py} (pipelines y rejillas), \texttt{train.py} (validación cruzada y reentreno),
\texttt{evaluate.py} (métricas en val y test) y \texttt{stacking.py}. Los parámetros comunes
(semilla, rutas, ventana, proporciones) están en \texttt{config.py}. Cada paso escribe en disco
y se salta si su salida ya está completa, lo que permite reanudar una ejecución cortada en el
CESGA. El Parquet limpio está particionado por \texttt{split}, de modo que cada lectura solo
carga los ficheros de la partición que necesita.

% ------------------------------------------------------------------------------------------
\section{Resultados}

\begin{table}
\centering
\caption{Resultados en test (177.060 textos, 22,5\,\% de IA). La línea base de la AUC-PR es
0,225. TPR@1\,\%: detección de IA con un 1\,\% de humanos marcados.}
\label{tab:resultados}
\begin{tabular}{llccccc}
\toprule
Features & Modelo & AUC-PR & AUC-ROC & F1 & FPR & TPR@1\,\% \\
\midrule
Léxicas & LR & 0,837 & 0,946 & 0,773 & 0,086 & 0,365 \\
Léxicas & SVC & 0,789 & 0,923 & 0,731 & 0,099 & 0,300 \\
Léxicas & NB & 0,692 & 0,889 & 0,651 & 0,149 & 0,197 \\
Léxicas sin \emph{stopwords} & SVC & 0,696 & 0,883 & 0,660 & 0,144 & 0,194 \\
Léxicas sin \emph{stopwords} & LR & 0,672 & 0,887 & 0,655 & 0,163 & 0,154 \\
Léxicas sin \emph{stopwords} & NB & 0,648 & 0,883 & 0,644 & 0,177 & 0,130 \\
Caracteres & NB & 0,594 & 0,855 & 0,587 & 0,274 & 0,115 \\
Estilo & GBT & 0,645 & 0,871 & 0,629 & 0,180 & 0,144 \\
Estilo & RF & 0,634 & 0,864 & 0,619 & 0,173 & 0,137 \\
Estilo & LR & 0,531 & 0,818 & 0,567 & 0,225 & 0,067 \\
Previsibilidad & GBT & 0,369 & 0,668 & 0,423 & 0,438 & 0,046 \\
Previsibilidad & RF & 0,368 & 0,668 & 0,423 & 0,424 & 0,046 \\
Previsibilidad & LR & 0,333 & 0,625 & 0,396 & 0,454 & 0,037 \\
\midrule
Stacking & LR & \textbf{0,863} & \textbf{0,957} & \textbf{0,796} & \textbf{0,080} & \textbf{0,408} \\
\bottomrule
\end{tabular}
\end{table}

La Tabla~\ref{tab:resultados} resume los resultados. Las features léxicas son las mejores, y
dentro de ellas la regresión logística. Quitar las \emph{stopwords} empeora tanto el SVC (de
0,789 a 0,696) como Naive Bayes (de 0,692 a 0,648): buena parte de la señal está en cómo se usan las palabras
funcionales, no en el vocabulario temático. En estilo, GBT supera claramente a la regresión
logística, así que las interacciones entre features importan. La previsibilidad es la familia
más débil: el modelo de trigramas, entrenado sobre todo con relatos de Reddit, mide más el
parecido con ese género que la previsibilidad del texto. Aun así, su regresión logística
detecta el 82\,\% de GPT-3, frente al 32\,\% de la de estilo, lo que apunta a información
complementaria (aunque en test solo hay 722 textos de GPT-3).

En todos los modelos val y test difieren menos de 0,01 en AUC-PR, así que el umbral elegido en
val se transfiere bien. La AUC-PR de la validación cruzada no es comparable con estas cifras:
train tiene un 29\,\% de IA y la AUC-PR depende de la prevalencia.

% ------------------------------------------------------------------------------------------
\section{Conclusiones}

Con features clásicas y Spark MLlib se detecta texto generado con una AUC-PR de 0,837, muy por
encima de la línea base, en un pipeline que procesa 1,4 millones de textos sin salir de la JVM.
Para llegar a un resultado creíble fue tan importante limpiar los atajos del dataset como
elegir el modelo, y el análisis por fuente muestra que aún queda uno, el género del texto.

\paragraph{Limitaciones.} Las fuentes son inferidas a partir del \texttt{id}, y el límite entre
IvyPanda y WebText mezcla unos 500 ensayos. Descartar la primera frase elimina algo de señal
real. Varios hiperparámetros elegidos están en el borde de la rejilla (\emph{smoothing} de
Naive Bayes, profundidad de los GBT), así que las rejillas podrían ampliarse.

\paragraph{Trabajo futuro.} Evaluar dejando fuera un generador o un dominio entero, para medir
la generalización a modelos no vistos; equilibrar los géneros entre clases; y sustituir el
modelo de trigramas por la perplejidad de un modelo neuronal pequeño.

\bibliographystyle{splncs04}
\bibliography{referencias}

\end{document}
